% !TEX program = lualatex
% Transparencias de Fundamentos de las Matemáticas I
% Clase 12: Operaciones internas

\documentclass[aspectratio=169,11pt]{beamer}

\usepackage{fontspec}
\usepackage[spanish,es-nodecimaldot]{babel}
\usepackage{amsmath,amssymb,mathtools}
\usepackage{booktabs}
\usepackage{csquotes}
\usepackage{xcolor}

\usetheme{Madrid}
\usecolortheme{default}
\setbeamertemplate{navigation symbols}{}
\setbeamertemplate{footline}[frame number]

\definecolor{FdMBlue}{RGB}{20,55,100}
\definecolor{FdMRed}{RGB}{130,30,30}
\definecolor{FdMGreen}{RGB}{25,100,60}

\setbeamercolor{structure}{fg=FdMBlue}
\setbeamercolor{title}{fg=white,bg=FdMBlue}
\setbeamercolor{frametitle}{fg=white,bg=FdMBlue}
\setbeamercolor{block title}{fg=white,bg=FdMBlue}
\setbeamercolor{block body}{bg=blue!3}
\setbeamercolor{alerted text}{fg=FdMRed}

\setmainfont{Latin Modern Roman}
\setsansfont{Latin Modern Sans}
\setmonofont{Latin Modern Mono}

\newcommand{\N}{\mathbb{N}}
\newcommand{\Z}{\mathbb{Z}}
\newcommand{\Q}{\mathbb{Q}}
\newcommand{\R}{\mathbb{R}}
\newcommand{\C}{\mathbb{C}}
\newcommand{\Pow}{\mathcal{P}}
\newcommand{\id}{\operatorname{id}}
\newcommand{\mcd}{\operatorname{mcd}}
\newcommand{\mcm}{\operatorname{mcm}}

\title[FdM I -- Clase 12]{Operaciones internas}
\subtitle{Asociatividad, conmutatividad, neutro e inverso}
\author{Antonio Falcó Montesinos}
\institute{Universidad CEU Cardenal Herrera}
\date{Fundamentos de las Matemáticas I}

\begin{document}

\begin{frame}
  \titlepage
\end{frame}

\begin{frame}{Objetivos de la clase}
\begin{itemize}
  \item Definir operación interna.
  \item Reconocer asociatividad y conmutatividad.
  \item Estudiar elementos neutros e inversos.
\end{itemize}
\end{frame}

\begin{frame}{Mapa de la clase}
\tableofcontents
\end{frame}

\section{Operaciones}

\begin{frame}{Operaciones}
\begin{block}{Operación interna}
\begin{itemize}
  \item Una operación interna en \(A\) es una aplicación \(*:A\times A\to A\).
  \item A cada par \((a,b)\) le asigna un elemento \(a*b\) de \(A\).
  \item La clausura dentro de \(A\) es parte de la definición.
\end{itemize}
\end{block}
\end{frame}

\begin{frame}{Operaciones}
\begin{block}{Ejemplos y contraejemplos}
\begin{itemize}
  \item La suma usual es interna en \(\mathbb{N}\).
  \item La resta usual no es interna en \(\mathbb{N}\), porque \(2-5\notin\mathbb{N}\).
  \item Una misma regla puede ser interna en un conjunto y no en otro.
\end{itemize}
\end{block}
\end{frame}

\section{Propiedades}

\begin{frame}{Propiedades}
\begin{block}{Asociatividad y conmutatividad}
\begin{itemize}
  \item Asociativa: \((a*b)*c=a*(b*c)\).
  \item Conmutativa: \(a*b=b*a\).
  \item Son propiedades de una operación concreta en un conjunto concreto.
\end{itemize}
\end{block}
\end{frame}

\begin{frame}{Propiedades}
\begin{block}{Elemento neutro}
\begin{itemize}
  \item Un elemento \(e\in A\) es neutro si \(e*a=a*e=a\) para todo \(a\in A\).
  \item Si existe, el elemento neutro es único.
  \item La unicidad se demuestra comparando dos posibles neutros.
\end{itemize}
\end{block}
\end{frame}

\begin{frame}{Propiedades}
\begin{block}{Elemento inverso}
\begin{itemize}
  \item Si hay neutro \(e\), un inverso de \(a\) es un elemento \(b\) con \(a*b=b*a=e\).
  \item La existencia de inversos depende del conjunto y de la operación.
  \item Ejemplo: \(2\) no tiene inverso multiplicativo en \(\mathbb{Z}\).
\end{itemize}
\end{block}
\end{frame}

\section{Profundización}

\begin{frame}{Profundización}
\begin{block}{Antes de estudiar propiedades}
\begin{itemize}
  \item Primero hay que comprobar que la operación está definida para todo par de elementos.
  \item Después se verifica que el resultado pertenece al mismo conjunto.
  \item Solo entonces tiene sentido preguntar por asociatividad, neutro o inversos.
\end{itemize}
\end{block}
\end{frame}

\begin{frame}{Profundización}
\begin{block}{Unicidad del neutro}
\begin{itemize}
  \item Supongamos que \(e\) y \(e'\) son neutros.
  \item Como \(e\) es neutro, \(e*e'=e'\).
  \item Como \(e'\) es neutro, \(e*e'=e\).
  \item Por tanto, \(e=e'\).
\end{itemize}
\end{block}
\end{frame}

\begin{frame}{Profundización}
\begin{block}{Dependencia del conjunto}
\begin{itemize}
  \item La suma tiene neutro en \(\mathbb{N}\), \(\mathbb{Z}\), \(\mathbb{Q}\) y \(\mathbb{R}\).
  \item Los inversos aditivos existen en \(\mathbb{Z}\), pero no todos existen en \(\mathbb{N}\).
  \item La estructura depende de la operación y del conjunto simultáneamente.
\end{itemize}
\end{block}
\end{frame}

\begin{frame}{Cierre}
\begin{itemize}
  \item Una operación interna es una función desde \(A\times A\) hasta \(A\).
  \item Las propiedades algebraicas se comprueban con cuantificadores.
  \item La próxima clase organiza estas propiedades en estructuras.
\end{itemize}
\end{frame}

\begin{frame}{Trabajo recomendado}
\begin{itemize}
  \item Releer las secciones correspondientes del manual.
  \item Identificar definiciones, símbolos nuevos y enunciados principales.
  \item Preparar dudas y ejemplos para el seminario asociado.
\end{itemize}
\end{frame}

\end{document}
